\section{Data Appendix}
\label{app:empirical}

\subsection{Sources}
All trip data come from the NYC Taxi and Limousine Commission's monthly
yellow-taxi parquet files \parencite{tlc2026trip}, which cover the
universe of metered yellow-cab trips. Taxi-zone geometry comes from the
TLC's published shapefile. No restricted or proprietary data are used;
the entire analysis is reproducible from public sources with the code in
this repository.

\subsection{Cell construction}
Two builds produce the estimation inputs. The first streams each monthly
file and aggregates to date-by-minute cells within windows around the
three clock cutoffs (14:30--17:30, 18:30--21:30, and 05:00--07:00),
recording for each cell the trip count, card-trip count, mean metered
surcharge, sum of card tips, sum of card pre-tip totals, the count of
tips exactly at a 20/25/30 percent button, and the zero-tip count. The
second aggregates by month, pickup zone, and \$5 fare bin over the
congestion-fee windows (2018m7--2019m12 and 2024m1--2025m6), recording
counts, tip and base sums, and mean fee incidence. Four further
builders complete the inputs: the 6:00am window cells (the main minute
build collects only the afternoon and evening windows; the 6:00am
builder is schema-identical), the JFK flat-fare minute cells
(RatecodeID~2, all three windows, same schema), the 2012--13 vendor
base-ladder cells, and the June-2025 tip-rate histogram cells behind
Figure~\ref{fig:spikes}.

Sample restrictions are applied identically in both builds: metered
fares between \$3 and \$200, non-negative tips. Tips are recorded only
for credit-card payments, so all tip means are conditional on card use;
the card share is reported as a locked outcome at every cutoff.
Date-by-minute cells are restricted to dates within each file's own
month, which removes a small number of misdated records.

\subsection{Regime and holiday coding}
The rush-hour surcharge rose from \$1.00 to \$2.50 and the overnight
surcharge from \$0.50 to \$1.00 on December 19, 2022
\parencite{tlc2022notice}; cells are assigned to regimes on that date.
Federal and municipal holidays, on which the rush surcharge does not
apply, are excluded from weekday samples and are listed in the cleaning
code.

\subsection{Geographic treatment assignment}
The congestion-fee columns exist only after each policy takes effect, so
a treated indicator built from them would place every pre-period
observation in the control group. We therefore learn the treated zone
set from post-period fee incidence---pickup zones where the mean fee
exceeds \$0.30, giving 52 zones for the 2025 CBD toll and 115 for the
2019 surcharge---and apply that fixed zone set to the whole window.
